\chapter{设计主动推断模型的配方}

\section{引言}

本章给出一个由四个步骤组成的配方，用于构建主动推断模型。我们将讨论为了实现一个模型必须作出的关键设计选择，并为这些选择提供若干指导原则。本章也为本书第二部分作铺垫；在那里，我们将展示若干基于主动推断的具体计算模型，以及它们在不同认知领域中的应用。

由于主动推断是一种规范性方法，它试图尽可能从第一性原理出发，解释行为、认知与神经过程。因此，主动推断的设计哲学是一种自上而下的哲学。与许多其他计算神经科学方法不同，这里的挑战并不是把大脑一块一块地仿真出来，而是找到那个能够刻画大脑试图解决之问题的生成模型。一旦这个问题被恰当地形式化为生成模型，在主动推断之下，问题的解就会自然显现，同时也会给出关于大脑与心智的相应预测。换言之，生成模型为所研究系统提供了一种完整描述。由此产生的行为、推断以及神经动力学，都可以通过最小化自由能从模型中推导出来。

生成建模方法已被广泛用于认知模型实现、统计建模、实验数据分析和机器学习等多个领域（Hinton 2007b; Lee and Wagenmakers 2014; Pezzulo, Rigoli, and Friston 2015; Allen et al. 2019; Foster 2019）。在这里，我们主要关心的是如何设计能够产生目标认知过程的生成模型。前几章已经展示过这种设计方法。例如，在预测编码的生成模型中，知觉被表述为对感觉最可能原因的推断；在离散时间演化的生成模型中，规划被表述为对最可能行动路线的推断。根据问题本身的不同（例如空间导航中的规划，或视觉搜索中的扫视规划），我们可以调整这些生成模型的形式，为它们赋予不同的结构（例如浅层或层级式）和变量（例如关于以环境为中心或以自我为中心的空间位置的信念）。关键在于，在对被优化生成模型的形式作出不同假设时，主动推断会呈现出多种不同面貌。例如，关于模型是在离散时间还是连续时间中演化的假设，会影响消息传递的形式（见第 4 章）。这意味着，对生成模型的选择，也对应着对行为和神经生物学的具体预测。

这种灵活性很有价值，因为它允许我们用同一种语言来描述多个领域中的过程。然而，从实践角度看，它也会带来困惑，因为为了找到研究对象的恰当描述层次，研究者必须作出一系列选择。在本书第二部分中，我们将通过一系列主动推断的计算示例来化解这种困惑。本章首先提出一个用于设计主动推断模型的一般配方，突出后续章节数值分析中会反复出现的一些关键设计选择、区分和二分法。

\section{设计主动推断模型：一个四步配方}

设计一个主动推断模型需要四个基础步骤，每一步都在回答一个特定的设计问题：

\begin{enumerate}
\item 我们在为哪个系统建模？首先必须确定感兴趣的系统。这一点并不像看上去那样简单，因为它依赖于对该系统边界，也就是马尔可夫毯的识别。什么算作主动推断主体（生成模型）？什么算作外部环境（生成过程）？它们之间的接口，也就是感觉数据与动作，又是什么？
\item 生成模型最合适的形式是什么？接下来的三个实践挑战中的第一个，是判断一个过程更适合被理解为类别性的（离散）推断，还是连续推断，这会推动我们在主动推断的离散时间实现、连续时间实现或混合实现之间作出选择。随后，我们需要选择最合适的层级深度，即在浅层模型与深层模型之间作出取舍。最后，我们还要考虑是否有必要赋予生成模型时间深度，以及预测依赖动作的观测的能力，以支持规划。
\item 如何搭建生成模型？生成模型最合适的变量和先验是什么？哪些部分是固定的，哪些部分必须通过学习获得？我们将强调选择正确变量和先验信念的重要性；同时也会强调时间尺度上的分离，即推断过程中较快更新的是状态变量，而学习过程中较慢更新的是模型参数。
\item 如何搭建生成过程？生成过程由哪些要素构成？它们与生成模型有何不同？
\end{enumerate}

在大多数情况下，这四个步骤就足以设计出一个主动推断模型。一旦完成，系统的行为便由主动推断的标准机制所决定：主动状态和内部状态沿着与模型相关的自由能泛函下降。从更实际的角度说，一旦我们明确了生成模型与生成过程，就可以使用标准的主动推断软件程序来获得数值结果，并完成数据可视化、分析和拟合（例如基于模型的数据分析）。下面我们将依次回顾这四项设计选择。

\section{我们在为哪个系统建模？}

把主动推断形式主义用于具体问题时，一个有用的第一步，是识别感兴趣系统的边界，因为我们关心的是：系统内部的东西如何通过感觉受体和效应器（例如肌肉或腺体）与外部世界发生交互。正如第 3 章所讨论的，形式化地区分系统内部状态、外部变量以及介于两者之间并中介其相互作用的变量，一种方法是诉诸马尔可夫毯（Pearl 1988）。再次强调，马尔可夫毯可以分为两类变量（Friston 2013）：一类中介外部世界对系统内部状态的影响，也就是感觉状态；另一类中介系统内部状态对外部世界的影响，也就是主动状态。见图 6.1。

需要强调的是，内部与外部之间的边界可以有多种划定方式。在本书第二部分将讨论的大多数模拟中，都会在一个主体（大致上是一个生物体）与其环境之间设置一种马尔可夫毯式的分离。这与认知模型的常规设定一致：一个主体依靠其内部状态（例如脑状态）实现知觉与动作选择等认知过程，并配备传感器和效应器。

但这并不是唯一可能的设定。从神经生物学角度看，我们可以把马尔可夫毯画在单个神经元周围，也可以画在大脑周围，或者画在整个身体周围。在第一种情况下，感觉状态包括突触后受体占据情况，主动状态则包括装载神经递质的囊泡与突触前膜融合的速率。神经元的内部状态（例如膜电位、钙离子浓度）于是就可以被看作是在依据某种隐含的生成模型，对其感觉状态的原因进行推断（Palacios, Isomura et al. 2019）。在这种设定里，被建模的外部状态包括该神经元所参与的神经网络。这与把整个网络都视为马尔可夫毯内部时发生的推断非常不同。比如，如果我们考察一个系统，它的感觉状态是视网膜中的光感受器，主动状态是眼动肌，那么内部状态所做的推断就是关于大脑之外的事物。这说明尺度十分重要，因为在这个马尔可夫毯下的内部状态，也包含了从单个神经元视角看见的内部状态。后一种内部状态在马尔可夫毯围绕单个神经元时，看起来是在推断大脑内部的事物；但当马尔可夫毯围绕整个神经系统时，情况就不是这样了。

这一点对于处理具身认知或延展认知视角尤其重要（Clark and Chalmers 1998; Barsalou 2008; Pezzulo, Lw et al. 2011）。例如，如果我们把马尔可夫毯画在神经系统周围，那么身体其他部分就会成为一种外部状态，我们必须通过内感受感觉状态对其作出推断（Allen et al. 2019）。或者，我们也可以把马尔可夫毯画在整个有机体周围。这样一来，看起来就像是除了大脑之外的其他器官也在对其环境进行推断。例如，皮肤在外部压力下发生凹陷，可以被表述为对外部压力来源的一种推断。延展认知观点更进一步，认为身体外部的对象也可以被纳入马尔可夫毯之内；例如，使用计算器来辅助推断，就意味着这个计算器成为推断系统内部状态空间的一部分。最后，我们还可以设想多个彼此嵌套的马尔可夫毯，例如大脑、有机体与共同体。

总之，定义马尔可夫毯能够确保我们知道：什么是被推断的对象（外部状态），以及什么在执行推断。事实上，相对于某个生成模型的自由能最小化，只涉及系统的内部状态和主动状态；这些状态只能够接触感觉状态，因此它们只能通过感觉状态间接地推断世界的外部状态。

\section{生成模型最合适的形式是什么？}

一旦我们确定了系统的内部状态，以及中介其与外部世界相互作用的状态，接下来就需要明确那个用来解释外部状态如何影响感觉状态的生成模型。正如前几章所述，主动推断可以作用于不同种类的生成模型。因此，我们必须为手头的问题指定最合适的生成模型形式。这主要意味着三项设计选择。第一，是在包含连续变量、离散变量或两者兼有的模型之间作出选择。第二，是在浅层模型与层级式或深层模型之间作出选择。浅层模型中的推断只在单一时间尺度上进行，也就是说所有变量都在同一时间尺度上演化；而层级模型中的推断则跨越多个时间尺度进行，不同变量会在不同时间尺度上演化。第三，是在只考虑当前观测的模型，与具有某种时间深度、能够考虑动作或计划后果的模型之间作出选择。

\subsection{离散变量还是连续变量，还是两者兼有？}

第一项设计选择，是考虑使用离散变量还是连续变量的生成模型更合适。前者包括对象身份、备选行动计划，以及连续变量的离散化表征。这些通常通过表达某个变量在每一个时间步上转移为另一种类型的概率来建模。后者则包括位置、速度、肌肉长度、亮度等变量，它们需要用变化率形式来表达生成模型。

从计算角度看，这一区分并不总是泾渭分明，因为连续变量可以被离散化，而离散变量也可以通过连续变量来表达。然而，在概念层面上，这一区分很重要，因为它支撑着关于目标认知过程时间进程是离散还是连续的具体假设。在当前大多数主动推断实现中，高层决策过程，例如在不同行动路线之间作出选择，通常用离散变量建模；而更细粒度的知觉和动作动力学，则通常用连续变量实现。我们将在第 7 章和第 8 章分别给出两类例子。

此外，离散变量与连续变量之间的选择，对神经生物学也有意义。虽然两种建模风格都诉诸自由能最小化，但它们所隐含的消息传递形式不同。如果人们认为消息传递与一种过程理论相关（见第 5 章），那么这意味着，实现自由能最小化的神经动力学在两类模型下也不相同。连续方案支撑的是预测编码，这是一种依赖自上而下预测并由自下而上预测误差进行校正的神经处理理论。而离散推断的类似过程理论，则涉及另一种形式的消息。最后，这两类模型也可以结合起来，使离散状态与连续变量相联系。也就是说，我们可以指定一个生成模型，其中某个离散状态（例如物体身份）会生成一组连续变量模式（例如亮度）。第 8 章将讨论一个同时包含离散变量和连续变量的混合生成模型示例。

\subsection{推断的时间尺度：浅层模型与层级模型}

第二项设计选择涉及主动推断所处的时间尺度。我们可以选择浅层生成模型，其中所有变量都在同一时间尺度上演化；也可以选择层级式或深层模型，其中包含在不同时间尺度上演化的变量，高层更慢，低层更快。

许多简单认知模型只需要浅层模型，但当目标认知过程中不同方面之间存在明确的时间尺度分离时，浅层模型就不够用了。语言加工就是一个例子：短时长的音素序列要由所说出的单词来提供语境，而短时长的单词序列又要由当前句子来提供语境。关键在于，一个单词的持续时间超出了组成它的任何单个音素，而一个句子的持续时间又超出了构成它的任何单个单词。因此，要为语言加工建模，我们可以考虑一种层级模型：句子、单词和音素分别位于从高到低的不同层级上，并以从慢到快、彼此近似独立的时间尺度演化。当然，这种分离只是近似的，因为不同层级之间必须相互影响；例如，句子会影响序列中的下一个单词，单词会影响序列中的下一个音素。

不过，这并不意味着为了开发针对某个单一层级的有意义模拟，我们就必须尝试建模整个大脑。例如，如果我们想聚焦于单词加工，那么就可以在不处理音素加工细节的情况下讨论它的某些方面。这意味着，从对音素进行推断的大脑区域传来的输入，可以被视为从单词加工区域角度看到的观测。用马尔可夫毯的语言来说，这通常意味着我们把模型较低层所执行的推断，当作该毯感觉状态的一部分。这样一来，我们就可以在不指定较低层、较快速推断过程细节的情况下，总结与当前时间尺度相关的推断；而这种层级分解也带来了巨大的计算收益。

另一个例子来自有意图的动作选择。在这种场景下，同一个目标，例如“进入你的公寓”，可以在较长时间内保持激活状态，并为一系列以更快时间尺度得到解决的子目标与动作提供语境，例如“找钥匙”“开门”“进去”。无论在连续域还是离散域中，这种时间尺度的分离都要求使用层级式、也就是深层生成模型。在神经科学中，我们可以假设皮层层级嵌入了这种时间尺度分离，即高层状态缓慢演化，低层状态快速演化，而这又重现了环境动力学的多时间尺度结构，例如在语音识别或阅读等知觉任务中所面对的环境动力学。在心理学中，这类模型有助于复现层级目标加工（Pezzulo, Rigoli, and Friston 2018）以及工作记忆任务（Parr and Friston 2017c），尤其是那些依赖延迟期活动的任务（Funahashi et al. 1989）。

\subsection{推断的时间深度与规划}

第三项设计选择涉及推断的时间深度。这里需要区分两类生成模型：第一类具有时间深度，会显式表征动作或动作序列，也就是策略或计划的后果；第二类则没有时间深度，只考虑当前观测而不考虑未来观测。图 4.3 展示了这两类模型的例子：上方的动态 POMDP，以及下方的连续时间模型。这两类模型的关键差别并不在于它们分别使用离散变量和连续变量，而在于只有前者，也就是具有时间深度的模型，赋予主体前瞻规划并在多个可能未来之间作出选择的能力。

想象一只啮齿动物在迷宫中规划通往已知食物位置的路线。要做到这一点，它需要一个具有时间深度的模型，这个模型大致等价于一种空间地图或认知地图（Tolman 1948），能够编码在动作条件下当前地点与未来地点之间的依赖关系，例如右转之后会到哪里，左转之后又会到哪里。动物可以利用这个具有时间深度的模型，反事实地考虑多条可能的行动路线，例如一系列左转和右转，并选择那条预期能够到达食物位置的路线。

为什么规划需要一个具有时间深度的模型？在主动推断中，规划是通过计算不同动作或策略所对应的期望自由能来实现的，然后选择那个具有最低期望自由能的策略。期望自由能并不仅仅是当前观测的函数，像变分自由能那样；它还是未来观测的泛函。未来观测按定义无法被直接观察，只能借助一个具有时间深度的模型来预测，而这个模型描述了动作如何产生未来观测。

因此，在设计主动推断主体时，一个有用的问题是：它是否应该拥有规划和面向未来的能力；如果答案是肯定的，就应该选择一个具有时间深度的模型。此外，还要考虑规划深度，也就是规划过程能够看多远的未来。最后，我们还可以设计既具有层级结构又具有时间深度的生成模型，使规划在多个时间尺度上展开，低层更快，高层更慢。是否对依赖策略选择的备选未来进行建模，在很大程度上也与离散模型和连续模型之间的选择纠缠在一起，因为“在多个由动作序列定义的备选未来之间进行选择”这一想法，用离散时间模型来表达通常更加直接。

\section{如何搭建生成模型？}

当我们已经明确研究的系统，并识别出生成模型的相关形式，例如连续或离散表征、浅层或层级结构之后，接下来的挑战，就是指定要纳入生成模型的具体变量，并决定其中哪些变量保持固定，哪些变量会因学习而发生变化。

\subsection{设置生成模型中的变量}

生成模型中的变量，既可以预先定义，也可以从数据中学习得到。为了说明问题，本书讨论的大多数模型都使用了预定义变量。在设计这类模型时，实践中的主要挑战是：判断哪些隐状态、观测和动作最适合当前问题。例如，第 2 章中那个能够区分青蛙和苹果的知觉模型，只包含两个隐状态（青蛙、苹果）和两个观测（跳、不跳）。一个更复杂的模型还可以加入额外观测，例如红色、绿色，以及诸如触摸之类的动作；这些动作在面对青蛙或苹果时，会产生不同的感觉结果，例如跳或不跳。

图 6.2 以示意方式展示了一个关于“会跳的青蛙”这一概念的生成模型。这个概念被表述为一个层级模型：图中央的一个单一的、多模态或超模态隐状态，会级联展开为对应不同模态知觉的一组单模态隐状态，最终在这些模态上引起感觉。这里的模态包括外感受、本体感受和内感受，见框 6.1。这样的安排相当于把“会跳的青蛙”这一概念视为多种感觉后果的共同原因，例如在视觉域中看到一个绿色且会跳的东西，在听觉域中听到蛙鸣；其中有些后果还可能依赖动作，例如当你去触碰它时，看到某物跳动的可能性会增加。对这一生成模型的反演，就对应于一种知觉推断：从它已被观察到的感觉后果，例如看到一个绿色而且会跳的东西，推断出“有一只会跳的青蛙”；并且这一推断会整合来自多个模态的信息。

\paragraph{框 6.1 感觉模态的类型：外感受、本体感受与内感受}

在主动推断中，通常会对三类感觉模态作概念区分：外感受模态，例如视觉和听觉；本体感受模态，例如对关节和肢体位置的感觉；以及内感受模态，例如对身体内部器官状态的感觉，如心脏和胃。在多模态生成模型中，往往可以把与不同模态相关的模型部分进行因子化；这样便能够表示，例如扫视运动会产生视觉后果，但不会产生听觉后果。

需要强调的是，主动推断的相同原则适用于所有这些模态。例如，正如视觉处理可以被描述为对知觉场景中某些隐变量的推断，内感受处理也可以被描述为对报告身体内部状态之隐变量的推断。进一步地，那些改变知觉场景的运动动作，以及那些改变内感受状态的内向动作，也可以用相似的方式来描述。前者会调动满足本体感受预测的脊髓反射，后者则会调动满足内感受预测的自主神经反射。这样的内感受处理支撑着异稳态与适应性调节，而其失调则可能带来精神病理学后果（Pezzulo 2013; Seth 2013; Pezzulo and Levin 2015; Seth and Friston 2016; Allen et al. 2019）。

一旦这些感兴趣的变量被确定下来，下一步工作就是把完整的生成模型写出来。图 2.1 中那个区分青蛙和苹果的简单生成模型就是一个例子：它完全由关于隐状态的先验信念，以及隐状态与观测之间的似然映射所指定，而这些数值既可以手工设定，也可以从数据中学习得到（见 6.5.2）。

超出这个简单例子之后，需要明确哪些元素，将完全由所选生成模型的形式决定。例如，图 4.3 上方展示的离散时间 POMDP 模型，需要指定 $A$、$B$、$C$、$D$ 和 $E$ 矩阵；连续方案也有与之类似的要素，只是命名没有这么按字母排序，第 8 章会讨论它们。但即使在这些更复杂的情形中，工作内容与上面的思路并没有本质差别：也就是为感兴趣的变量指定先验信念，例如在离散时间实现中，对第一个时间步的隐状态使用 $D$ 向量，对观测使用 $C$ 矩阵；并指定它们之间的概率映射，例如在 $A$ 矩阵中给出隐状态与观测之间的似然映射。不过，在某些情况下，从生成模型状态空间的因子化角度来思考会很有帮助，因为如果某些变量并不必要，就可以避免考虑所有可能的变量组合。在第 7 章中，我们会讨论一个发生在知觉处理中的、具有生物学合理性的因子化例子，即“是什么”通路与“在哪里”通路之间的分工（Ungerleider and Haxby 1994）；也就是分别表示物体身份和位置的变量，它们在模型中往往可以被独立处理，从而简化模型，因为二者经常彼此不变。

决定哪些变量值得纳入模型，以及它们在模型中如何相互关联或被因子化，往往是模型设计中最困难、但也最具创造性的部分。它本质上是在把我们的认知假设翻译成一种支持主动推断的数学形式。我们应当如何选择“正确”的变量？归根到底，这个问题可以表述为：提出若干看起来合理的备选模型，然后选择自由能最低的那个，这与贝叶斯模型比较的精神一致。不过，从大多数研究的实践角度看，一个很有用的观点是：生成模型应当尽可能接近我们认为数据是如何被生成出来的方式。当我们在认知心理学场景中诉诸主动推断时，这通常意味着，我们应当去想：实验心理学家是如何生成并呈现给被试的刺激的。当这些过程被形式化为所需的概率分布之后，我们就得到了一个生成模型，而对其自由能的最小化动力学，会自然导向相关任务的完成。

这里可以类比大多数关于知觉的贝叶斯模型或理想观察者模型：这些模型被设计成在很大程度上模仿当前任务的结构，就像识别青蛙还是苹果的例子那样（第 2 章）。这一想法有时被称为良好调节器定理（Conant and Ashby 1970），其含义是：为了有效调节环境，一个生物体，无论是生物性的还是人工的，都必须是该系统的良好模型。从生态位建构的角度看，这有时又会被表述为：生物体模型与环境之间具有某种统计适配性（Bruineberg et al. 2018），反之亦然。不过，这并不意味着一个主体的生成模型必须与真实生成数据的生成过程完全相同。对于大多数实践应用来说，它可以更简单，也可以有所不同。稍后在本章的 6.6 节中，我们还会回到这一点。

\paragraph{框 6.2 先验与经验行为}

关于如何选择先验，还有另一种视角，它来自一组被称为完全类定理的结果（Wald 1947; Daunizeau et al. 2010）。这些结果指出，任何统计决策程序，也就是任何行为，都可以在一组合适的先验信念之下被表述为贝叶斯最优。这意味着，如果我们想解释经验行为，挑战就在于识别出那样一个生成模型，也就是一组先验信念，使得该行为能够尽可能简单地被重现。简言之，先验表达的是关于目标系统的一种假设。如果还存在其他同样合理的先验信念，这就为借助经验数据进行贝叶斯模型比较提供了机会。

这一点对临床人群中的计算表型分析也有含义。既然总会存在某组先验信念使行为在贝叶斯意义下成为最优，那么，在理解那些产生精神病学或神经学综合征的计算缺陷时，关键问题就不再是行为是否最优，而是这些先验究竟是什么。这个想法起初可能有些违反直觉。然而，完全类定理意味着，单纯问一个行为是不是贝叶斯最优，本身并没有意义。真正重要的问题是：什么样的先验信念会使这种行为成为最优？在第 9 章中，我们将看到，借助以科学家自身信念为基础的自由能最小化，我们可以回答这个问题。

\subsection{生成模型的哪些部分是固定的，哪些部分是学习得到的？}

另一项设计选择，是决定生成模型中哪些部分保持固定，哪些部分会随着时间推移并由于学习而被更新。原则上，主动推断允许模型的每一个部分，甚至包括模型结构本身，随着时间而被更新或学习。因此，学习是一项设计选择，而不是某种必需项。相应地，我们既会讨论完全由手工设计出来的主动推断模型，也会讨论这样一些模型：其中某些部分，例如转移概率，保持固定，而另一些部分，例如似然，则会随着时间而被更新。

在主动推断中，学习被表述为推断的一个方面，也是一种自由能最小化过程。到目前为止，我们把推断描述为对生成模型状态之信念的更新。以完全类似的方式，我们也可以把学习描述为对生成模型参数之信念的更新。为此，生成模型必须被赋予关于待学习分布参数的先验信念，而这些具体参数取决于每个变量所对应的概率分布形式，例如高斯分布中的均值和方差。每当遇到新数据时，这些先验值就会被更新，形成后验信念。正如我们将在第 7 章讨论的，这种更新在算法形式上与状态变量的更新是相同的。

推断和学习都依赖同一种贝叶斯式信念更新，这一点在模型设计时可能会引起混淆；部分原因是，什么应该被建模为状态，什么应该被建模为参数，并不总是显而易见。然而，在认知模型中，推断和学习之间仍然存在明确差别。推断描述的是我们关于模型状态之信念的快速变化，例如当我们看到一个红色物体后，如何更新“面前是一个苹果”的信念。学习描述的是我们关于模型参数之信念的缓慢变化，例如在多次观察到红苹果之后，如何更新似然分布，从而提高“苹果--红色”映射的数值。关于参数的信念通常比关于状态的信念变化得慢得多，而且它们往往只有在状态已被推断之后才会更新。从神经生物学角度看，把推断映射为神经动力学、把学习映射为突触可塑性，是很有吸引力的。进一步地，正如我们将在第 7 章看到的，对模型参数持有概率性信念会诱发新奇寻求行为，使主体选择最有利于学习其世界因果结构的数据。这说明，赋予主动推断模型以学习其参数，甚至其结构的能力，是研究主动学习与基于好奇心探索之行为动力学的一种有效方式。

在结束本节之前，值得指出的是：本书所举的例子大多是相对简单的生成模型，它们使用表格式方法来定义，例如用显式矩阵表示先验和似然，并且运行在较小的状态空间中。相比之下，在机器学习、深度学习和机器人学等领域，人们正在发展更复杂的生成模型及其相应的学习方案，例如变分自编码器（Kingma and Welling 2014）、生成对抗网络（Goodfellow et al. 2014）、递归皮层网络（George et al. 2017）以及世界模型（Ha and Schmidhuber 2018）。原则上，我们可以借用其中任何一种方法，或者更多其他方法，来实现主动推断模型的一个或多个组成部分，例如似然模型或转移模型。通过利用最新的机器学习方法，主动推断有望被扩展到越来越具有挑战性的领域和应用；例如可参见 Ueltzhöffer（2018）和 Millidge（2019）。

不过，在设计使用复杂机器学习模型的主动推断模型时，有几点必须认真考虑，尤其是在我们关心其认知和神经生物学含义的时候。主动推断的吸引力之一，在于它为认知功能提供了一种整合性视角，认为知觉推断、动作规划和学习等都来自同一个自由能最小化过程。如果我们把若干彼此独立运行或学习的生成模型简单并列，这种整合力量就会丧失。此外，上述机器学习方法本身对应于与主动推断不同的过程模型，因此也会带来不同的认知和神经生物学解释。最后，当使用机器学习方法时，这里讨论的某些设计选择，例如关于模型变量的选择，可能会被跳过，因为它们成为学习中涌现出来的属性；但与此同时，它们又会被另一类设计选择取代，例如深度神经网络的层数、参数数量和学习率等。这些设计选择同样可能具有重要的认知和神经生物学含义，但已经超出了本章的讨论范围。

\section{搭建生成过程}

在主动推断中，生成过程描述的是主动推断主体之外的世界动力学，也就是决定主体观测的那个过程（见图 6.1）。看起来似乎有些奇怪：我们直到介绍完主体的生成模型之后，才来定义生成过程。毕竟，一个建模者通常从一开始就会心里有一个任务，也就相应地有一个生成过程；因此，把顺序倒过来，先设计生成过程，再设计生成模型，似乎完全说得通。特别是在生成模型需要通过情境中的交互学习获得的应用里，例如游戏式环境或机器人场景（Ueltzhöffer 2018; Millidge 2019; Sancaktar et al. 2020），这种顺序看起来尤其自然。

我们之所以把生成过程的设计推迟到这里，是因为在本书讨论的许多实践应用中，我们只是简单地假定：生成过程的动力学与生成模型相同，或者非常相似。换句话说，我们通常假设主体的生成模型会紧密模仿生成其观测的过程。这并不等于说主体对环境拥有完美知识。事实上，即便主体知道是什么过程生成了它的观测，它也仍然可能对该过程中的某些因素不确定，例如对其初始状态不确定，正如苹果与青蛙的例子那样。用离散时间主动推断的语言来说，我们完全可以设计这样一个模型：生成模型与生成过程具有相同的 $A$ 矩阵，但主体关于初始状态的信念，即作为其生成模型一部分的 $D$ 向量，却与生成过程的真实初始状态不同，甚至彼此不一致。这里有一个微妙之处需要注意：即使生成模型和生成过程都由相同的 $A$ 矩阵和 $B$ 矩阵来刻画，它们的语义也并不一样。生成过程中的 $A$ 矩阵是环境的一种客观属性，在贝叶斯模型中有时称为测量分布；而生成模型中的 $A$ 矩阵则编码了主体的主观信念，在贝叶斯模型中称为似然函数。

当然，除了最简单的情形之外，生成模型与生成过程并不必须相同。在主动推断的实践实现中，我们总是可以把生成过程与生成模型分开指定：既可以使用不同于生成模型的方程，也可以使用其他方法，例如游戏模拟器，这类系统以动作作为输入、以观测作为输出（Cullen et al. 2018），从而遵循图 6.1 所隐含的、由马尔可夫毯界定的通常行动--知觉回路。

设计与生成过程相似或不相似的生成模型，还具有一些哲学含义（Hohwy 2013; Clark 2015; Pezzulo, Donnarumma et al. 2017; Nave et al. 2020; Tschantz et al. 2020）。如前所述，良好调节器定理（Conant and Ashby 1970）认为，一个有效的适应性生物体必须拥有，或者本身就是，它所调节系统的一个良好模型。不过，这可以通过多种方式实现。第一种方式，也是到目前为止我们一直在讨论的方式，是让生物体的生成模型去模仿生成过程，至少在很大程度上如此。以这种方式构建的模型，可以称为显式模型或环境模型，因为它们的内部状态与环境的外部状态之间存在相似性。第二种方式，则是让生物体的生成模型比生成过程更加简约，甚至与生成过程显著不同，只要它能够正确处理环境中那些对适应性行动和目标达成真正有用的方面。以这种方式构建的模型，可以称为感觉运动模型或行动导向模型，因为它们主要编码的是动作--观测，或者说感觉运动依赖关系，其首要作用是支持目标导向行动，而不是提供对环境的精确描述。

如果我们考虑对同一个场景的不同建模方式，这一区别会更容易理解。例如，一只啮齿动物试图逃离一个迷宫，而迷宫中的某些通道是死胡同。一个显式生成模型可能会像一张迷宫的认知地图，并对位置、走廊和死胡同等外部实体进行详细刻画。借助这种模型，啮齿动物可以依靠基于地图的导航走出迷宫。而一个行动导向模型，则可能编码胡须动作与触觉之间的依赖关系。这样的模型能够支持情境上合适的策略选择，例如当没有体验到或预期到触觉时就向前移动，而在相反情况下改变方向。最终，啮齿动物也能够借此逃出迷宫，而无需显式表征位置、走廊或死胡同。这两类模型会引出对主动推断不同的哲学解释：生成模型究竟是在重建外部环境，还是在提供精确的动作控制条件。

最后，正如形态计算领域所讨论的那样（Pfeifer and Bongard 2006），生物体或机器人的控制有些方面可以外包给身体本身，因此并不需要被编码进其生成模型。一个例子是被动动力步行器：这种物理对象在外形上类似人体，由两条“腿”和两只“臂”组成，却能够在没有传感器、电机或控制器的情况下沿斜坡行走（Collins et al. 2016）。这个例子意味着，至少某些运动能力方面的功能，可以通过精心调谐的身体力学来实现，以便利用环境依赖关系，例如合适的体重或体型使个体能够在不打滑的情况下行走；因此，这些依赖关系不必被编码进生物体的生成模型。这提示了设计主动推断主体及其身体的另一种方式：不是让它们拥有环境的良好模型，而是让它们本身就成为环境的良好模型。不过，这些不同的主动推断模型设计方式并不是互斥的，而是可以根据具体问题恰当地组合起来。

\section{使用主动推断进行模拟、可视化、分析与数据拟合}

在大多数实践应用中，一旦生成模型和生成过程都被定义出来，我们只需要使用主动推断的标准程序，也就是让主动状态和内部状态沿着与模型相关的自由能泛函下降，就可以得到数值结果。可以说，建模者的目标通常是模拟、可视化、分析以及拟合数据，例如开展基于模型的数据分析。支持这些功能的标准主动推断程序可以免费获得（\texttt{https://www.fil.ion.ucl.ac.uk/spm/}）；附录 C 给出了一个使用这些程序的注释示例。

尽管在大多数情况下，主动推断程序都可以开箱即用，但在某些实践应用中，人们仍然可能考虑做一些特定的微调或改动。例如，指定规划的时间深度，会决定在计算期望自由能时要考虑多少未来状态。对于大型状态空间中的主动推断实践应用来说，设置一个有限的时间深度，再结合其他替代穷举搜索的近似方法，例如采样（Fountas et al. 2020），可能会很有帮助。

另一个对主动推断标准运行方式进行调整的例子，是有选择地移除期望自由能方程中的某些部分。这种消融操作有助于把标准主动推断，即使用完整期望自由能的系统，与若干简化版本进行比较；在这些简化版本中，期望自由能的某些部分被压制，从而使其在形式上与 KL 控制或效用最大化系统等框架变得可比（Friston, Rigoli et al. 2015）。此外，我们也可以为主动推断模型增加一些额外机制，例如习惯学习（Friston, FitzGerald et al. 2016）或学习率调制（Sales et al. 2019）；不过需要注意的是，如果想保持主动推断的规范性特征，就必须把这些附加机制同样表述为自由能最小化过程。

最后，对主动推断作出其他一些微调或改变，也有助于刻画推断障碍和精神病理状态。例如，我们可以通过赋予某个生物体的生成模型过强或过弱的先验，也就是过高或过低水平的神经调质，来探索其行为和神经后果。第 9 章将给出若干与精神病理学相关的主动推断模型示例。

\section{总结}

本章概述了在搭建主动推断模型时必须作出的最重要设计选择。我们提供了一个四步配方，并给出了一些指导原则，用来处理模型设计者通常会面对的挑战。当然，并不需要以僵硬的方式遵循这套配方。有些步骤可以颠倒顺序，例如先设计生成过程再设计生成模型，也可以把若干步骤合并。但总体上说，这些步骤都是必需的。

这也为本书余下内容奠定了基础。接下来的章节将通过一系列示例把这些思想付诸实践，以展示本书前半部分提出的理论原则。后文所有示例之间真正不同的地方，只在于我们在这里强调过的这些设计选择。第二部分将展示具有不同边界的系统，展示离散或连续动力学在不同时间尺度上的实现，也将展示在许多不同领域中，先验信念的选择对于重现行为为何至关重要；但它们都在实现同一个主动推断框架。
